\frontchapter{Abstract}

Last-mile delivery platforms increasingly buy capacity from a contracted
fixed fleet with a known price, from gigworkers who drive open multi-stop
routes, and from occasional drivers who carry parcels on trips they make
anyway. Crowd drivers' value of time is private, and paying what they claim
invites overstatement. This thesis studies a procurement mechanism in which
truthful reporting is a dominant strategy for both crowd classes: routes are
generated from public data before bidding, winners are selected exactly, and
Vickrey--Clarke--Groves (VCG) payments are paid.

Because the route pool must be fixed before bids are seen, its size is the
bottleneck. A gigworker has no detour budget to bound the search, so her
pool grows as $\Theta(n^B)$ in the number of orders $n$ for a bundle cap $B$,
and only 0.4--1.5\% of generated routes are ever optimal on the test
instances. The thesis asks how far such a pool can be pruned by a rule that
reads neither bids nor other drivers' data.

Routes are generated by a forward label-setting algorithm. An envelope that
completes subset routes with the publicly priced fixed fleet defines a
\emph{local frontier}. The thesis proves that deleting every route outside
it leaves all payments unchanged, because every deleted route has a
substitute that is itself retained. Results planned for the final thesis
concern the minimality and worst-case size of the frontier, a label rule
that prunes during enumeration, and a decomposition of the payment
computation.

The methods were implemented in Python with CPLEX and evaluated in a
pre-registered study on 1,500 synthetic instances with 10 to 20 orders;
route generation was also run on up to 100 orders, and the full evaluation
will be extended to that scale. All solves were optimal and all checks against a brute-force oracle passed. The
frontier removed 75--100\% of the pool without changing any payment, and
where the fixed fleet dominates it certified before bidding that 60\% of the
drivers with a feasible route need not take part. Joint bidding, bundling
and truthful payments mattered only where the crowd competes with the fixed
fleet: there they lowered system cost by 5.8\% and 11.3\% and carried a
14.1\% payout premium (medians). The fixed-fleet price decides which regime
applies. Future work should study rules that use other drivers' public data
and richer driver reports.

\medskip
\noindent\textbf{Keywords:} Crowdshipping; Mechanism design; VCG auction;
Pickup and delivery; Column pruning
